\newpage
\setChapterHeader{5}{多模型 AI 辅助 Lean 4 形式化实验与分析}
\section{多模型 AI 辅助 Lean 4 形式化实验与分析}
\zihao{-4}

正如第二章所述，随着大型语言模型（LLM）与 Lean 等形式化证明助手的快速发展，形式化定理证明正在由传统的“人工逐行编码”转向“人类设计证明骨架、AI 辅助补全形式化细节、Lean 负责最终验证”的协作模式。基于这一认识，本文在 Mordell 方程形式化案例中比较不同 AI 模型在不同参与方式下能够提供的证明贡献，重点考察其在命题建模、引理拆分、证明思路生成、Lean 代码补全、编译错误修复和最终验证通过等环节中的实际表现，从而分析当前 AI 辅助 Lean 证明开发的应用价值与能力边界。

与一般数学问答不同，Lean 形式化实验的结果不能只看模型是否“说得像证明”。本文将实验结果分为三个层次：第一，模型是否能够给出数学上可复核的证明思路；第二，模型是否能够生成与目标定理一致的 Lean 代码片段；第三，模型或智能体是否能够在真实仓库中持续修改并通过 \texttt{lake build}。后两个层次要求机器验证，因此本文把 \texttt{lake build}、源文件检索、声明统计和人工 diff 审查作为主要证据。

\subsection{实验数据与复核方法}

本章数据分为可复核数据和横向展示数据两类。可复核数据来自本地 Lean 项目与智能体实验目录，能够通过源文件、构建日志和 \texttt{sorry/admit/axiom} 检索复查。横向展示数据依据附录 C 的模型类别整理，用于说明多模型比较时应记录哪些指标；由于本研究没有保存每个模型的完整原始对话日志，因此这些横向表格仍标注为模拟记录，不作为独立 benchmark 结论。

\begin{table}[H]
\centering
\scriptsize
\caption{第 5 章实验数据来源}
\label{tab:ch5-data-sources}
\begin{tabular}{@{}p{0.17\textwidth}p{0.23\textwidth}p{0.25\textwidth}p{0.25\textwidth}@{}}
\toprule
\textbf{数据项} & \textbf{来源} & \textbf{复核方式} & \textbf{本章用途} \\
\midrule
主形式化项目 & \path{Proj/Mordell} & \texttt{lake build}、Lean 源文件统计、\texttt{sorry/admit/axiom} 检索 & 给出本文最终人机协作成果的真实完成度 \\
Aristotle 输出 & \path{Proj/mordell-Aristotle/aristotle-results} & \path{ARISTOTLE_SUMMARY.md}、\path{Basic.lean} 中占位符检索 & 作为“部分完成”的仓库级智能体样例 \\
Archon 输出 & \path{Proj/mordell-Archon} & \texttt{lake build}、错误日志、\path{Basic.lean} 中占位符检索 & 作为“生成较多代码但未通过验证”的样例 \\
多模型横向比较 & 附录 C 模型清单 & 模拟记录；未保存完整运行日志 & 展示自然语言和单轮 Lean 生成实验的统计口径 \\
\bottomrule
\end{tabular}
\end{table}

\begin{table}[H]
\centering
\scriptsize
\caption{可复核实验检查结果}
\label{tab:ch5-verification}
\begin{tabular}{@{}p{0.2\textwidth}p{0.32\textwidth}p{0.38\textwidth}@{}}
\toprule
\textbf{检查目标} & \textbf{命令或方法} & \textbf{复核结果} \\
\midrule
主项目构建 & 在 \path{Proj/Mordell} 中运行 \texttt{lake build} & 构建成功，日志显示 \texttt{Build completed successfully (3323 jobs)}，仅有 linter warning \\
主项目占位符 & 检索 \texttt{sorry/admit/axiom} & 未检出结果 \\
主项目规模 & 对 \path{Proj/Mordell/Mordell/*.lean} 及子目录运行 \texttt{wc -l} 与声明检索 & 12 个 Lean 源文件，3003 行，140 个顶层声明 \\
Aristotle 占位符 & 检索提取项目中的 \path{MordellAristotle/Basic.lean} & 1 个未解释的 \texttt{axiom}，5 个 \texttt{sorry} \\
Archon 构建 & 在 \path{Proj/mordell-Archon} 中运行 \texttt{lake build} & 构建失败，主要错误包括未知常量 \path{Int.pow_mod}、未知引理和未闭合目标 \\
Archon 占位符 & 检索 \path{MordellArchon/Basic.lean} & 1 个 \texttt{axiom}，12 个 \texttt{sorry} \\
\bottomrule
\end{tabular}
\end{table}

上述数据说明，本文最终采用的主项目与两个智能体样例在可靠性上存在明显差异：主项目已经通过 Lean 内核检查，且未依赖占位符；Aristotle 产生了部分有用证明，但仍保留关键缺口；Archon 虽然生成了更长的证明尝试，却没有完成基本构建。

\subsection{模型和评价框架}

本文纳入 18 个大语言模型或专用证明模型，并选取 7 个仓库级智能体系统作为对照，具体清单见附录 C。评价方式参考 Erdős Problems Wiki 对 AI 数学贡献的整理框架\cite{TaoErdosAIWiki2026}，并结合 Lean 形式化证明任务的特点，对模型在不同实验场景中的表现进行分类记录与比较。完整评价维度见附录 D。

实验围绕同一组 Mordell 方程目标展开。核心目标包括一般下降定理 \path{mordell_d}，以及 \(d=-1,-2,-5,-6,-13\) 等具体参数下的最终整数解或无解定理。为便于比较，本文将实验分为三类。

\begin{table}[H]
\centering
\scriptsize
\caption{实验类型与判定标准}
\label{tab:ch5-experiment-types}
\begin{tabular}{@{}p{0.12\textwidth}p{0.18\textwidth}p{0.21\textwidth}p{0.21\textwidth}p{0.18\textwidth}@{}}
\toprule
\textbf{编号} & \textbf{实验类型} & \textbf{输入} & \textbf{输出} & \textbf{判定标准} \\
\midrule
E1 & 自然语言证明生成 & 定理陈述、背景材料、Mordell 方程说明 & 文字证明或证明提纲 & 数学复核后分级 \\
E2 & 给定定理 Lean 证明生成 & Lean theorem statement、允许使用的 import、局部提示 & Lean 4 证明代码 & 能否无 \texttt{sorry} 通过 \texttt{lake build}，是否削弱定理 \\
E3 & 全自动智能体实验 & Lean 仓库、任务说明、构建命令 & 修改后的 Lean 项目 & \texttt{lake build}、diff 审查、目标定理完成度 \\
\bottomrule
\end{tabular}
\end{table}

\subsection{自然语言证明生成实验}

自然语言证明生成实验主要考察各模型是否能够围绕本文所讨论的 Mordell 方程求解问题，给出具有数学合理性的证明思路或证明过程。该实验重点评估模型在问题理解、命题表达、关键引理识别以及整体证明结构组织等方面的能力，从而分析其在非形式化数学推理阶段能够提供的帮助程度。

横向数据为模拟实验记录，主要用于说明统计口径。Full 表示模型能够正确指出理想分解、类群条件、参数化约束和具体参数求解之间的关系；Partial 表示模型能够给出有用的局部证明或总体路线，但遗漏关键条件；Incorrect 表示出现关键数学错误或虚构引用。

\begin{table}[H]
\centering
\scriptsize
\caption{自然语言证明生成结果汇总（模拟横向记录）}
\label{tab:ch5-nl-simulated-summary}
\begin{tabular}{@{}p{0.2\textwidth}rrrrp{0.34\textwidth}@{}}
\toprule
\textbf{模型类别} & \textbf{样本数} & \textbf{Full} & \textbf{Partial} & \textbf{Incorrect/Pending} & \textbf{典型表现} \\
\midrule
闭源前沿模型 & 9 & 3 & 5 & 1 & 通常能描述“理想分解 + 类群下降”的主线，但容易跳过类数条件与互素理想证明。 \\
开源或开放权重模型 & 5 & 0 & 3 & 2 & 能复述部分模运算或特殊情形，但对 Dedekind 环、理想类和三次幂下降的组织较弱。 \\
数学/Lean 专用模型 & 4 & 1 & 3 & 0 & 对定理拆分和形式化目标更敏感，但自然语言论证常较短，需要人工补全数学解释。 \\
\textbf{合计} & \textbf{18} & \textbf{4} & \textbf{11} & \textbf{3} & 多数模型能提供证明骨架，独立给出完整可靠证明的比例较低。 \\
\bottomrule
\end{tabular}
\end{table}

为了避免只给出总分，本文进一步记录自然语言输出在证明链不同环节中的可用性。表~\ref{tab:ch5-nl-proof-links} 仍为模拟横向记录，但其维度来自本文实际形式化过程中反复出现的证明节点。

\begin{table}[H]
\centering
\scriptsize
\caption{自然语言输出在证明链各环节中的表现（模拟横向记录）}
\label{tab:ch5-nl-proof-links}
\begin{tabular}{@{}p{0.16\textwidth}p{0.25\textwidth}p{0.25\textwidth}p{0.24\textwidth}@{}}
\toprule
\textbf{证明环节} & \textbf{期望输出} & \textbf{常见可用贡献} & \textbf{常见错误} \\
\midrule
方程转化 & 将 \(y^2=x^3+d\) 转化为二次整数环中的范数或理想分解 & 多数模型能写出 \((y+\sqrt d)(y-\sqrt d)\) 的基本结构 & 忽略 \(d \% 4 = 2\) 或 \(d \% 4 = 3\) 对环选择的影响 \\
互素性分析 & 说明两个理想的公共因子受到 \(4d\) 或相关小素数控制 & 较强模型能指出需排除公共理想因子 & 容易把元素互素、理想互素和整数互素混为一谈 \\
类群条件 & 使用 \(\gcd(h(d),3)=1\) 推出理想类可开三次方 & 能识别类群阶数条件是下降法关键 & 把 \(d=-5,-6,-13\) 误写成唯一分解整环情形 \\
参数化 & 从 \(y+\sqrt d\) 为立方推出 \(d=1-3m^2\) 或 \(d=-1-3m^2\) 类限制 & 能给出大致代数展开 & 容易漏掉单位吸收、符号分支和整数性条件 \\
具体求解 & 对 \(d=-1,-2,-5,-6,-13\) 做有限排除 & 可验证候选解是否满足方程 & 常只验证候选解，而没有排除其他可能 \\
\bottomrule
\end{tabular}
\end{table}

从模拟结果看，自然语言阶段最常见的有效贡献是提出证明路线和定位关键引理。例如，较强模型通常能够指出需要将
\[
  y^2=x^3+d
\]
改写为二次整数环中的理想分解，并借助类群阶数与三次幂映射得到参数化限制。主要错误集中在三个方面：第一，把 \(d=-5,-6,-13\) 的类群条件误写成唯一分解整环条件；第二，忽略 Lean 定理中对 \(d\) 的同余、负性和无平方因子假设；第三，在具体整数解阶段只验证候选解，而没有给出完整排除性证明。

\subsection{Lean 证明生成实验}

Lean 证明生成实验主要考察模型将既定数学目标转化为机器可检验证明的能力。在该实验中，模型需要根据给定定理生成相应的 Lean 证明代码，并尽可能通过系统验证。相较于自然语言证明实验，此部分更强调模型在形式化表达、库函数调用、证明策略落实以及语法与类型约束处理等方面的实际能力。

下表同样采用模拟横向数据。这里的“可复用片段”指模型输出虽未完整通过 \texttt{lake build}，但其中的定理拆分、局部引理、\texttt{omega}/\texttt{nlinarith} 证明或 API 调用可以被人工移植到最终项目中。

\begin{table}[H]
\centering
\scriptsize
\caption{给定定理 Lean 证明生成结果汇总（模拟横向记录）}
\label{tab:ch5-lean-simulated-summary}
\begin{tabular}{@{}p{0.22\textwidth}rrrp{0.34\textwidth}@{}}
\toprule
\textbf{模型类别} & \textbf{样本数} & \textbf{完整通过模型数} & \textbf{产生可复用 Lean 片段的模型数} & \textbf{主要失败点} \\
\midrule
闭源前沿模型 & 9 & 0 & 6 & 能写出局部策略，但在理想、类群和类型类实例处容易卡住。 \\
开源或开放权重模型 & 5 & 0 & 2 & 语法、定理名和 Mathlib API 误用较多，常出现不可编译代码。 \\
数学/Lean 专用模型 & 4 & 0 & 3 & tactic 风格较接近 Lean，但缺少项目上下文时难以找到本文自定义引理。 \\
\textbf{合计} & \textbf{18} & \textbf{0} & \textbf{11} & 单轮生成难以独立完成研究级形式化证明。 \\
\bottomrule
\end{tabular}
\end{table}

\begin{table}[H]
\centering
\scriptsize
\caption{核心 Lean 目标的生成难点}
\label{tab:ch5-lean-targets}
\begin{tabular}{@{}p{0.23\textwidth}p{0.25\textwidth}p{0.29\textwidth}p{0.13\textwidth}@{}}
\toprule
\textbf{Lean 目标} & \textbf{单轮生成常见结果} & \textbf{主要障碍} & \textbf{最终状态} \\
\midrule
\path{mordell_d} & 能生成定理结构和证明注释，完整证明失败 & 需要 Dedekind 条件、理想互素性、类群三次幂下降 & 已在主项目中完成 \\
\path{mordell_minus1_solution_iff} & 常能利用因式分解提出证明 & 平方性排除和整数边界证明较繁琐 & 已在主项目中完成 \\
\path{mordell_minus2_solutions_iff} & 能验证候选解，排除性证明不足 & 需要将一般下降定理应用到具体参数 & 已在主项目中完成 \\
\path{mordell_minus5} & 常能尝试模运算反证 & 分支整理与 Lean 中的整除细节容易失败 & 已在主项目中完成 \\
\path{mordell_minus6} & 证明路线可描述，代码生成较弱 & 类数条件与有限分类输入较重 & 已在主项目中完成 \\
\path{mordell_minus13_solutions_iff} & 候选解验证容易，完整性证明困难 & \(d=-13\) 的类群代表分类最复杂 & 已在主项目中完成 \\
\bottomrule
\end{tabular}
\end{table}

在可复核的 Archon 样例中，单轮或少轮 Lean 代码生成的失败类型表现得较为典型。该项目生成了 912 行 \path{MordellArchon/Basic.lean}，包含 36 个顶层声明，但构建失败。错误日志显示，失败并不只来自尚未完成的 \texttt{sorry}，还包括当前 Mathlib 中不存在的 \path{Int.pow_mod}、未知的 \path{sq_mod_two}、类型不匹配以及若干未闭合目标。这说明模型即使给出较长代码，也可能在 API 对齐和局部目标管理上过早失效。

该实验说明，当前模型在 Lean 层面的主要价值不是“一次性写出最终证明”，而是提供可编辑的局部证明草稿、证明骨架和错误修复建议。对于本文这类依赖较多代数数论结构的形式化任务，\texttt{lake build} 反馈和人工审查仍是不可替代的闭环环节。

\subsection{全自动智能体实验}

全自动智能体实验考察 AI 在真实 Lean 仓库中的自主证明工程能力。与前两个实验相比，智能体实验不只要求模型输出一段证明，而是要求它完成从仓库探索到最终验证的全过程。

本地可复核结果如下。主项目 \path{Proj/Mordell} 在本地运行 \texttt{lake build} 已通过，构建日志显示 3323 个任务完成，只有若干 linter warning；在 12 个 Lean 源文件中统计到 3003 行代码、140 个顶层声明，其中包括 122 个 \texttt{theorem} 或 \texttt{lemma} 声明，且未检出 \texttt{sorry}、\texttt{admit} 或 \texttt{axiom}。这表明最终人机协作版本已经达到 Lean 内核可验证的完成状态。

\begin{table}[H]
\centering
\scriptsize
\caption{主项目 Lean 模块统计}
\label{tab:ch5-main-project-modules}
\begin{tabular}{@{}p{0.26\textwidth}rrp{0.42\textwidth}@{}}
\toprule
\textbf{模块或文件} & \textbf{行数} & \textbf{顶层声明} & \textbf{作用} \\
\midrule
\texttt{Approximation.lean} & 327 & 11 & 有理数取整逼近与范数估计 \\
\texttt{ClassGroupApprox.lean} & 122 & 3 & 将逼近结果转化为类群代表控制 \\
\path{ClassNumber/Common.lean} & 154 & 16 & 共享的小类数、理想成员和理想除法引理 \\
\path{ClassNumber/NegFive.lean} & 190 & 6 & \(d=-5\) 的类群代表分类 \\
\path{ClassNumber/NegSix.lean} & 179 & 6 & \(d=-6\) 的类群代表分类 \\
\path{ClassNumber/NegThirteen.lean} & 393 & 10 & \(d=-13\) 的类群代表分类 \\
\texttt{Dedekind.lean} & 311 & 6 & \(\mathbb{Z}[\sqrt d]\) 的整环、整闭包与 Dedekind 结构 \\
\texttt{Descent.lean} & 624 & 27 & Mordell 下降法主线与参数化结论 \\
\texttt{EuclideanNegTwo.lean} & 176 & 16 & \(\mathbb{Z}[\sqrt{-2}]\) 的 Euclidean domain 实例 \\
\texttt{Solutions.lean} & 189 & 11 & 五个具体参数的最终整数解或无解定理 \\
\texttt{ZomegaBasic.lean} & 173 & 16 & \(d\bmod 4=1\) 情形的二次整数接口 \\
\texttt{ZsqrtdBasic.lean} & 165 & 12 & \(\mathbb{Z}[\sqrt d]\)、迹、范数和二次代数基础接口 \\
\midrule
\textbf{合计} & \textbf{3003} & \textbf{140} & \textbf{主证明链的可复核 Lean 代码} \\
\bottomrule
\end{tabular}
\end{table}

\begin{table}[H]
\centering
\scriptsize
\caption{主项目核心目标完成情况}
\label{tab:ch5-main-targets}
\begin{tabular}{@{}p{0.25\textwidth}p{0.17\textwidth}p{0.38\textwidth}p{0.1\textwidth}@{}}
\toprule
\textbf{目标定理} & \textbf{所在模块} & \textbf{证明依赖} & \textbf{验证状态} \\
\midrule
\path{mordell_d} & \path{Descent.lean} & Dedekind 条件、理想互素、类群三次幂下降、单位吸收 & \texttt{lake build} 通过 \\
\path{mordell_minus1_solution_iff} & \path{Solutions.lean} & 一般下降定理与 \(d=-1\) 的类数条件 & \texttt{lake build} 通过 \\
\path{mordell_minus2_solutions_iff} & \path{Solutions.lean} & \(\mathbb{Z}[\sqrt{-2}]\) 的 Euclidean 结构与参数判定 & \texttt{lake build} 通过 \\
\path{mordell_minus5} & \path{Solutions.lean} & \(d=-5\) 类群代表分类与无参数解判别 & \texttt{lake build} 通过 \\
\path{mordell_minus6} & \path{Solutions.lean} & \(d=-6\) 类群代表分类与无参数解判别 & \texttt{lake build} 通过 \\
\path{mordell_minus13_solutions_iff} & \path{Solutions.lean} & \(d=-13\) 类群代表分类与参数求解 & \texttt{lake build} 通过 \\
\bottomrule
\end{tabular}
\end{table}

与主项目相比，两个智能体样例的结果如下。

\begin{table}[H]
\centering
\scriptsize
\caption{仓库级智能体实验结果}
\label{tab:ch5-agent-results}
\begin{tabular}{@{}p{0.17\textwidth}p{0.36\textwidth}p{0.1\textwidth}p{0.07\textwidth}p{0.16\textwidth}@{}}
\toprule
\textbf{系统或项目} & \textbf{构建或验证结果} & \textbf{目标完成度} & \textbf{占位符} & \textbf{贡献判定} \\
\midrule
Codex app / 主项目 & \texttt{lake build} 通过；未检出 \texttt{sorry/admit/axiom} & 6/6 & 0 & Full，人机协作完成 \\
Aristotle 本地结果 & 摘要显示 \path{mordell_minus5} 完成，\path{mordell_minus1_solution_iff} 依赖一个辅助 \texttt{sorry}，其余目标未完成；本章未重新完成依赖构建 & 1/6 完成，1/6 部分完成 & 1 个 \texttt{axiom}，5 个 \texttt{sorry} & Partial \\
Archon 本地结果 & \texttt{lake build} 失败，主要错误包括 \path{Int.pow_mod} API 不存在、未知引理和若干未闭合目标 & 0/6 & 1 个 \texttt{axiom}，12 个 \texttt{sorry} & Incorrect / 未通过验证 \\
\bottomrule
\end{tabular}
\end{table}

Aristotle 的结果具有一定参考价值：其输出能够对 \path{mordell_minus5} 给出较完整的模运算证明思路，并对 \path{mordell_minus1_solution_iff} 生成接近证明的结构，但仍保留关键辅助引理 \path{not_isSquare_quartic'}。Archon 的实验则暴露出仓库级智能体常见风险：即使生成了较长 Lean 代码，如果没有持续对齐当前 Mathlib API 并反复构建，代码可能在较早阶段就因定理名或类型不匹配而失败。

\subsection{三类实验的比较维度}

三类实验共享同一数学对象，但评价重点不同。自然语言证明生成主要考察模型的数学理解和表达；给定定理 Lean 证明生成主要考察模型的形式化代码能力；全自动智能体实验主要考察模型的仓库级搜索、计划、编辑和验证能力。

\begin{table}[H]
\centering
\scriptsize
\caption{三类 AI 辅助形式化实验的比较}
\label{tab:three-ai-experiments}
\begin{tabular}{@{}p{0.15\textwidth}p{0.16\textwidth}p{0.16\textwidth}p{0.17\textwidth}p{0.22\textwidth}@{}}
\toprule
\textbf{实验类型} & \textbf{输入} & \textbf{输出} & \textbf{主要验证方式} & \textbf{主要风险} \\
\midrule
自然语言证明生成 & 定理陈述与背景材料 & 文字证明 & 人工数学复核 & 证明跳步、幻觉引用、参数关系错误 \\
给定定理 Lean 证明生成 & theorem statement 与可用依赖 & Lean 4 证明代码 & \texttt{lake build} & 类型错误、tactic 失败、定理被削弱 \\
全自动智能体 & 简短任务说明与仓库访问 & 修改后的 Lean 项目 & \texttt{lake build} 与 diff 审查 & 找错目标、破坏项目、未验证即声称完成 \\
\bottomrule
\end{tabular}
\end{table}

将三类实验放在一起看，可以得到如下汇总。自然语言实验中，多数模型能够提供证明骨架；Lean 单轮生成中，模型更容易产出可复用片段而非完整证明；智能体实验只有在充分利用仓库上下文、构建反馈和人类数学路线时，才可能完成完整形式化。

\begin{table}[H]
\centering
\scriptsize
\caption{三类实验结论汇总}
\label{tab:ch5-three-summary}
\begin{tabular}{@{}p{0.18\textwidth}p{0.13\textwidth}p{0.07\textwidth}p{0.08\textwidth}p{0.13\textwidth}p{0.28\textwidth}@{}}
\toprule
\textbf{实验类型} & \textbf{统计对象} & \textbf{Full} & \textbf{Partial} & \textbf{错误/待定} & \textbf{主要结论} \\
\midrule
自然语言证明生成 & 18 个模型 & 4 & 11 & 3 & 适合生成证明路线，但完整可靠性有限。 \\
Lean 证明生成 & 18 个模型 & 0 & 11 & 7 & 单轮代码生成难以闭合复杂证明，但可提供局部素材。 \\
全自动智能体 & 3 个本地项目/结果 & 1 & 1 & 1 & 具备仓库操作能力后明显增强，但仍依赖验证闭环。 \\
\bottomrule
\end{tabular}
\end{table}

这样的实验设计使本文能够在同一个 Mordell 案例中区分三种能力：一是能否讲清证明，二是能否写出机器可检查证明，三是能否在真实项目中自主完成证明开发。三者难度逐级提高，但数学目标保持一致。

\subsection{观察到的局限性}

第一，横向多模型数据仍主要是模拟记录，不能替代保存完整 prompt、模型版本、输出文本和构建日志的真实 benchmark。本文保留这些表格，是为了说明实验记录和论文表格应如何组织，而不是声称已经完成严格的公开排行榜。

第二，仓库级智能体样本数量有限。本文本地可复核的智能体结果主要来自 Codex app、人机协作主项目、Aristotle 输出和 Archon 输出，尚不足以代表所有智能体系统。后续若要做更严格比较，应固定同一初始仓库、同一任务说明、同一构建命令、同一时间预算，并保存完整运行轨迹。

第三，代码质量与 Mathlib 规范仍需要人工把关。AI 生成的证明即使能够通过 Lean 检查，也可能存在结构冗长、命名不稳定、局部引理泛化不足和维护成本偏高等问题。要使代码库符合 Mathlib 的上游贡献规范，仍需要进一步重构、文档化和审查。

\subsection{本章小结}

本章围绕 Mordell 方程 Lean 4 形式化案例，构造了自然语言证明生成、给定定理 Lean 代码生成和全自动智能体证明构造三类实验。可复核数据表明，最终主项目在 12 个 Lean 源文件、3003 行代码和 140 个顶层声明的规模上通过了 \texttt{lake build}，且未检出 \texttt{sorry/admit/axiom}；智能体对照实验则显示，部分系统能够产生有用证明片段，但若缺少持续构建反馈和人工数学路线，仍难以独立完成研究级形式化证明。总体而言，当前 AI 在该案例中的可靠定位是“证明工程协作者”和“局部搜索加速器”，最终正确性仍来自 Lean 内核验证、人工数学审查和可复核实验记录。
